\FloatBarrier\Needspace{11\baselineskip}\section{结论}

Qiushi Engine 在 BabyLM Strict-Small 上主导并完成了前沿模型构建、原理发现和原理指导改进三个完整研究过程，持续开展问题形成、方法设计、实验执行、机制分析与成果整合。三个阶段围绕同一个科学问题推进：有限文本怎样让模型学会利用上下文中的信息关系，并使新学习与已有语言功能共存。

研究辨明了两个重要条件。逐字重复与对齐重述选择性地塑造来源信息的使用方式，作用取决于目标关系和预测窗口；熟悉输入上的表现恢复，也不保证未见输入仍能使用已学计算。由此提出的设计原则，是围绕目标所需的上下文依赖组织经验，分别设计可见信息、监督目标与功能保持，再以新旧能力的共同表现检验效果。

这一原则进入第二代模型的具体训练，在共同母模型的两次续训中均取得优于普通续训的完整评测结果。两代公开代表模型的 Overall 从 42.02 提高到 42.25。预算再分配、残差增量学习、关系锚点、身份匹配和测量控制等研究，还留下了各有问题与实验依据的独立成果。

对长期自主科研而言，这项工作展示了 Research RSI 的一个具体实例。Qiushi Engine 从自身研究中积累科学认识、方法技术与实验经验，使后续研究能够提出更明确的问题、构造新的方法，并用实际结果检验和修正设计依据。在这一过程中，原有成果既被继承，也成为新的研究对象；有效方法得到发展，失效解释被修正，研究路线在分化与整合中向前推进。

两代模型、学习原理、实验方法与独立发现共同构成可继续使用的研究基础。报告说明它们为何形成、如何相互联系，配套工程保存模型、代码、数据构造、完整结果和科研笔记，使同行能够检验已有认识，并据此开展新的研究。
